% GENERATED, edit the French unit or the translation table
\mxcorrections{2}\label{C02-CORR}
\textbf{Teacher use.} The answers below concern the tasks not worked in the lesson. First have students state their method; give graded help before the solution. The durations are estimates to be calibrated in class.
\subsection*{Diagnostic and remediation}
\begin{enumerate}
\item $x^2-6x+9$. \textit{Action:} Rebuild the product of two binomials.
\item $(x-4)(x+4)$. \textit{Action:} Use the difference of two squares on two examples.
\item $x=2$ or $x=-5$. \textit{Action:} Have students say aloud the zero-product rule and its hypothesis.
\item $x=-3$ or $x=3$. \textit{Action:} Distinguish the positive square root from the solutions of the equation.
\item Negative before 3, zero at 3, positive after 3. \textit{Action:} Draw the line, then make the sign table.
\item 4, reached only at 1. \textit{Action:} Start from the fact that a square is non-negative.
\end{enumerate}
\subsection*{Investigation activity}
For width $x$, the length is $12-x$; the geometric domain is $\intoo{0}{12}$. The requested areas are respectively 20, 32, 36, 32 and 20 square units. For every admissible width, $A(x)=x(12-x)=36-(x-6)^2\leq36$, with equality only at $x=6$. Thus the maximum-area rectangle is a square of side 6. Trials suggest the result; the nonnegative square supplies the proof.
\subsection*{Immediate applications}
\begin{enumerate}
\item $g(x)=(x+1)^2-4=(x-1)(x+3)$. Roots 1 and $-3$, minimum $-4$ at $-1$.
\item The first is $3(x-1)^2=0$, so 1. For the second, $\Delta=-16<0$: no real solution.
\item $f(x)=a(x-1)(x-4)$; $4a=-8$, so $f(x)=-2(x-1)(x-4)$.
\item $(x-5)(x+1)<0$ gives $\intoo{-1}{5}$. The vertex of $x(8-x)$ is at 4, outside the domain; on $\intff{0}{3}$ the function is increasing and its maximum is 15 at 3.
\end{enumerate}
\subsection*{Exercises: method and results}
\paragraph{Exercise 2.1: Completing the square.} $x^2-8x+7=(x-4)^2-9=(x-1)(x-7)$. The leading coefficient is positive; the quadratic is negative or zero between the roots: $\intff{1}{7}$.
\paragraph{Exercise 2.2: Discussing without heavy calculation.} $x^2+4>0$ always, no root. $(x+2)^2\geq0$, double root $-2$. For the last one, roots $-4$ and 1; positive on $\intoo{-4}{1}$, negative outside, zero at the endpoints.
\paragraph{Exercise 2.3: Choosing a strategy.} Special product: $x=\pm3$. Sum and product: 4 and 5. For the third, no obvious integer root: $\Delta=17$ and $x=\dfrac{3\pm\sqrt{17}}{4}$. The discriminant works in every case, but is not always the most efficient.
\paragraph{Exercise 2.4: Recovering a parabola.} $f(x)=a(x+2)(x-3)$. At 1: $-6a=12$, so $a=-2$; $f(x)=-2x^2+2x+12$, hence $f(0)=12$.
\paragraph{Exercise 2.5: A vegetable garden against a wall.} The two widths use $2x$ metres; length $20-2x$. Domain $\intff{0}{10}$ for the function, $\intoo{0}{10}$ for a non-degenerate rectangle. $A(x)=20x-2x^2=50-2(x-5)^2$. Maximum 50 m$^2$, dimensions 5 m and 10 m.
\paragraph{Exercise 2.6: A parameter.} $\Delta=4-4m$. Two roots if $m<1$; double root 1 if $m=1$; none if $m>1$. The parameter m stays fixed when we solve for x.
\paragraph{Exercise 2.7: Link with sequences.} $u_n=(n-1)(n-4)$; negative for $1<n<4$. Since n is an integer, the only indices are 2 and 3. The set of indices is not the whole real interval.
\paragraph{Exercise 2.8: Delayed review.} The minimum is $-5$ at $-3$. The inequality gives $(x+3)^2\leq1$, that is $-1\leq x+3\leq1$, so $\intff{-4}{-2}$.
\paragraph{CQFD, exercise 43 p.~99, questions a and c.} a: obvious root 1, product $-30$, so the second root is $-30$. c: write $4x^2+9x+5$; root $-1$, product $\dfrac54$, so the second root is $-\dfrac54$. Check the sums, $-29$ and $-\dfrac94$ respectively.
\paragraph{CQFD, exercise 49 p.~99, questions a and b.} The numbers are the roots of $t^2-st+p$. a: $t^2-t+1$, discriminant $-3$: no real pair. b: $t^2+60t+756$, discriminant $576=24^2$: $-18$ and $-42$. Their sum is $-60$ and their product 756.
\subsection*{MCQ: answers and reasons}
\textbf{MCQ 1.}
\begin{enumerate}
\item C. Set each factor to zero.
\item A. $4-12=-8$.
\item D. The square vanishes at 2, value 5.
\item B. $-\dfrac{b}{a}=\dfrac{6}{2}=3$.
\item C. Positive leading coefficient, strict inequality.
\end{enumerate}
\textbf{MCQ 2.}
\begin{enumerate}
\item D. Special product.
\item A. $\dfrac{c}{a}=-\dfrac{2}{3}$.
\item C. The function increases up to 2; on this domain, right endpoint 1.
\item B. For a nonzero, a zero discriminant characterises the double root.
\item D. The multiplying coefficient is free and nonzero.
\end{enumerate}
\subsection*{True or false}
\begin{enumerate}
\item False: $x^2+1$ is strictly positive.
\item False without a hypothesis on a; $-(x-1)(x-2)$ is positive between 1 and 2.
\item True: the square of a real number is positive or zero; equality only at 2.
\end{enumerate}
\paragraph{Differentiated teaching.} For a technical block, offer only the first step (domain, factor or reference sample space), then withdraw the help. For an independent student, ask for a second method and the treatment of a limiting case. Orally, allow an initial reformulation in the language the student has mastered, then require a precise mathematical sentence in French.
